\section{Introduction}
Researchers adapting a predictor to a target domain often have several source-data pools to choose from. The pools may differ in acquisition conditions, image style, or the examples they contain. To compare combinations, one can train a predictor on each and evaluate it on target validation data. The expense comes from repeating this process, even when each predictor is small. We ask whether an earlier comparison can reduce the number of combinations that receive supervised evaluation without discarding useful choices.

Pretrained encoders provide a common space for this comparison. An encoder maps each image to a feature vector. Keeping its parameters fixed allows those features to be reused across source combinations and classifier fits; CLIP and DINOv2 make this setting natural \citep{radford2021,oquab2024}. A classifier trained on these fixed features is called a readout. Before fitting it, we can compare how source pools cover different feature directions and how that coverage relates to the target.

Broad coverage alone does not establish usefulness for the target. Effective rank summarizes how evenly variation is spread across directions by exponentiating the entropy of normalized singular values \citep{roy2007,garrido2023rankme}. A high value may reflect variation rarely encountered at prediction time, and individually diverse pools may overlap. Source diversity and target relevance are therefore different questions. A score that uses target features can distinguish directions that a source-only spectral summary treats alike.

Geometry still leaves uncertainty about prediction: feature coverage does not by itself determine the relation between features and labels. We therefore use it to retain alternatives, then use labeled validation to choose among them. This separates two questions. Does screening leave good candidates available, and does validation select a good predictor from that set? A useful ranking must ultimately be judged by the final choice and the work needed to obtain it.

We evaluate a two-stage data-pool screening protocol. Unlabeled source and target features first rank combinations. Predictors are then trained on the retained source combinations and selected using separate target-validation data. Target-weighted experimental design favors coverage that reduces uncertainty along target-relevant directions \citep{yu2006}. Second-moment matching instead compares the source's average directional energy with the target's. We evaluate these scores against target-blind diversity, separate screening omissions from validation errors, and measure final loss across evaluation budgets and computation costs.

On equal-sample-cost DomainNet combinations, short geometric lists achieve lower mean test Brier, a probability-prediction loss, than longer random lists. Independent timing shows savings with cached features; adding measured feature work preserves the logistic-regression saving but reverses it for ridge regression. Source-domain controls and fixed-parameter repartitioning distinguish useful block information from domain-level matching, with outcomes depending on construction and metric. Projection and readout changes retain high average recall, whereas PACS and Office-Home leave little room to improve on direct-selection baselines. Together, these comparisons assess what target-aware geometry contributes to supervised pool selection and when its extra work is worthwhile.

\section{Related Work}
\subsection{Optimal design and distribution matching}
A-optimal design measures uncertainty through the trace of a covariance matrix, with applications to active linear regression \citep{fontaine2021}. Transductive experimental design minimizes prediction variance over given target covariates \citep{yu2006}; its regularized linear objective is the target-weighted trace used here. We apply this criterion to unions of fixed source pools and retain several unions for labeled validation. Distribution matching offers a different approach. MMD compares mean embeddings of source and target distributions \citep{gretton2012}; a homogeneous quadratic kernel makes this comparison depend on their second moments. Moment alignment has also been used to learn representations for multi-source domain adaptation \citep{peng2019}. Here the encoder stays fixed, allowing us to compare uncertainty reduction and moment matching as scoring rules on the same candidate combinations.

\subsection{Selection under a computation budget}
Compute-constrained data selection accounts for selection work as well as subsequent training \citep{yin2025}. Pool screening introduces an intermediate decision: how many combinations should receive supervised evaluation? A broad shortlist may retain good candidates but still require substantial fitting and validation. We compare both retention and the final selected model, then separately measure decision time and feature-preparation costs. This distinguishes fewer evaluations from a faster decision.

\subsection{Representation geometry and diversity}
RankMe uses effective rank to assess learned representations without downstream labels \citep{garrido2023rankme}. Determinantal point processes favor subsets with complementary directions \citep{kulesza2012}. Both motivate ways to compare source pools before fitting a task-specific model. Their geometric objectives differ, however, from preserving candidates with low target loss: variation may be redundant for a particular task, and directions important to that task may occupy only part of the representation. We examine merged effective rank and a specified subspace-DPP implementation as target-blind comparators.

\section{Geometry-Guided Screening}
\label{sec:problem}
\label{sec:theory}
\subsection{Problem and two-stage procedure}
A pool is a fixed block of examples; a candidate is a feasible union of blocks. The study uses a candidate family and separate target data for screening, validation, and testing. Screening and supervised evaluation produce one selected predictor. A shortlist contains candidate combinations, not individual images.

Let $H_i\in\R^{n_i\times d}$ contain the features of pool $i$, using a fixed encoder and shared label-independent preprocessing. Its Gram matrix $G_i=H_i^\top H_i$ records feature-direction strengths and their relationships. For $S$ in the candidate family $\cC$, $G_S=\sum_{i\in S}G_i$ and $n_S=\sum_{i\in S}n_i$. Unlabeled target-selection features estimate the uncentered second moment $C_T=\E_T[hh^\top]$, which describes target directional energy.

The screen ranks $\cC$ and retains its first $s$ combinations as $\cQ_s$. We then train predictors on their source labels and use a separate target-validation set to choose $\widehat S\in\cQ_s$. Target-test data evaluate this frozen choice and the candidates retained. The screening stage, denoted U0, sees source and target-selection features but no labels or validation/test data. The supervised stage, U2, reads source labels and target-validation data. Target-blind rules use only source information. This separation lets the screen exploit reusable features while reserving task-specific evidence for the final choice.

Each combination receives the same classifier and validation procedure; screening changes which combinations are evaluated, not the encoder. The equal-sample-cost study fixes training sample count, not elapsed time.

\subsection{Shared geometry and spectral diversity}
The singular value decomposition $H=U\Sigma V^\top$ separates feature directions from their strengths. For fixed nonnegative weights $w_i$, stacking the scaled source matrices gives
\begin{equation}
 G_S=H_S^\top H_S=\sum_{i\in S}w_iG_i.
 \label{eq:gram}
\end{equation}
The eigenvalues of this sum are the squared singular values of the pooled features. It preserves the strength and cross-pool alignment of directions, provided preprocessing coordinates remain fixed. The experiments use unit weights.

Effective rank normalizes the singular values into a probability vector and exponentiates its entropy \citep{roy2007}:
\begin{equation}
 \reff(H)=\exp\!\left(-\sum_j p_j\log p_j\right),\qquad
 p_j=\frac{\sigma_j(H)}{\sum_k\sigma_k(H)}.
 \label{eq:rank}
\end{equation}
For a nonzero matrix, $0\log0=0$ and $1\leq\reff(H)\leq\operatorname{rank}(H)$. Equal positive singular values attain the upper endpoint. Two singular values in ratio $9:1$ give effective rank about $1.38$, compared with two for equal strengths. We take square roots of Gram eigenvalues before normalization; normalizing the eigenvalues themselves would define a different score. Repeating an entire matrix leaves effective rank unchanged, so it measures spectral shape rather than sample evidence. It is a target-blind comparator, not an intermediate step required by the target-aware scores.

Marginal spectra and principal angles alone need not determine even this pooled statistic: the same strengths can be assigned differently to shared and unshared directions. Appendix~\ref{app:proofs} supplies an explicit counterexample and the conditions for a local spectral formula. We retain Gram operators when scoring combinations so that this directional information is available.

\subsection{Target weighting and moment matching}
Target weighting asks which source combination reduces uncertainty where the target needs coverage. A shared linear model makes this question precise. Consider scalar Bayesian regression with $\theta\sim N(0,P_0)$, $P_0\succ0$, and $y=h^\top\theta+\epsilon$, where independent noise satisfies $\epsilon\sim N(0,\sigma^2)$ and $\sigma^2>0$. Sources and target share $\theta$ and the noise model; target covariates reveal no further information about $\theta$. Conditional on the source design, the posterior covariance and posterior-mean target risk are
\begin{equation}
 P_S=\left(P_0^{-1}+\sigma^{-2}G_S\right)^{-1},\qquad
 R_T(S)=\sigma^2+\tr(C_TP_S).
 \label{eq:risk}
\end{equation}
The expectation averages posterior uncertainty and a new target observation. The derivation is recalled in Appendix~\ref{app:proofs}.

Equation~\eqref{eq:risk} motivates the score $f_A(S)=\tr(\widehat C_TP_S)$, minimized with fixed hyperparameters. With an isotropic prior, this is the regularized target-prediction-variance criterion of transductive experimental design \citep{yu2006}; we retain the name Target A-opt for continuity with the recorded experiments. Source observations reduce uncertainty in the directions they cover, and the target moment weights those reductions by how often the directions occur at prediction time. Thus a pool can be valuable even without high overall diversity if it covers directions that remain uncertain and matter to the target. Adding a positive-semidefinite source Gram cannot increase this score.

For the experimental choice $P_0^{-1}=I$ and $\sigma^2=1$, write $G_S=U\operatorname{diag}(\lambda_j)U^\top$, including zero eigenvalues. Then
\begin{equation}
 f_A(S)=\sum_j\frac{u_j^\top\widehat C_Tu_j}{1+\lambda_j}.
 \label{eq:target-directions}
\end{equation}
The denominator reflects source coverage of direction $u_j$, while the numerator measures its importance to the target. Two combinations with the same eigenvalues can therefore receive different scores if their directions align differently with the target. Effective rank treats their spectra identically; target weighting distinguishes where their coverage is useful.

Second-moment matching asks a different question: how close is the source's average directional energy to the target's? It minimizes
\begin{equation}
 f_M(S)=\norm{G_S/n_S-\widehat C_T}_F^2.
 \label{eq:mmd}
\end{equation}
This is empirical squared MMD with kernel $k(h,h')=(h^\top h')^2$, including diagonal terms. It brings the source's average directional energy closer to the target's rather than weighting a posterior uncertainty matrix. Both criteria read the same features, so their comparison tests whether inverse-covariance weighting adds value beyond direct moment matching.

The shared conditional model supplies the risk interpretation of A-opt. Its assumption is substantive: if all features equal $1$, source pools labeled constantly $+1$ and $-1$ produce predictors whose risk ordering reverses between two target worlds with opposite constant labels. The feature geometry is unchanged. In classification, we therefore evaluate the scores using trained ridge-softmax and logistic-regression readouts rather than identify Equation~\eqref{eq:risk} with their losses.

\subsection{Separating retention from final selection}
The final choice can lose quality in two ways: screening discards a good candidate, or validation chooses a worse member of the retained set. Let $\ell(S)$ be target risk for a fixed trained predictor and $v(S)$ its validation loss. If $|v(S)-\ell(S)|\leq\varepsilon$ throughout $\cQ_s$, the validation minimizer satisfies
\begin{equation}
 \ell(\widehat S)-\min_{S\in\cC}\ell(S)
 \leq
 \underbrace{\min_{S\in\cQ_s}\ell(S)-\min_{S\in\cC}\ell(S)}_{\text{shortlist omission}}
 +2\varepsilon.
 \label{eq:decomposition}
\end{equation}
Appendix~\ref{app:proofs} proves the bound. Its first term measures loss from excluding the best candidate; the second bounds validation error within the shortlist. We audit finite-test counterparts and top-ten retention. The deviation condition explains the two error sources; the experiments do not estimate it as a population guarantee.

\subsection{Decision budgets and cost accounting}
The full-search reference trains all candidates and chooses $S_{\mathrm{full}}=\arg\min_{S\in\cC}v(S)$. This is a deployable procedure. The test oracle, in contrast, is the candidate with lowest observed test loss and is used only for auditing. Let $b(S)$ be test Brier, the mean squared distance between predicted class probabilities and the one-hot label; lower is better. We report $b(\widehat S_s)-b(S_{\mathrm{full}})$ and the fraction of repeated tasks exceeding a declared absolute tolerance. The gap can be negative because the validation and test rankings can differ.

With frozen features already available, a screening run costs its ranking work plus fitting and validating the retained candidates. Starting from images adds encoder initialization and feature extraction; target-aware scores additionally need target-selection features. These are different accounting scenarios. The break-even condition is that screening work be smaller than the supervised evaluations it replaces, after any additional feature work is included. Candidate counts measure one part of that condition, not elapsed time. In particular, the same matrix operations that make a linear score convenient can be a substantial part of fitting a ridge readout.

\section{Evaluation Protocol}
\label{sec:protocol}
Controlled synthetic tasks test the risk model. PACS and Office-Home \citep{li2017,venkateswara2017} then evaluate direct selection against a prespecified practical-effect threshold. An exhaustive headroom analysis follows those outcomes. The subsequent DomainNet study fixes sample cost and evaluates shortlists; its configuration was frozen before feature extraction and outcome evaluation. Test-sample resampling and the later projection/readout study are exploratory follow-ups. Appendix~\ref{app:protocol} records their configurations and result locations.

From the official cleaned DomainNet data \citep{peng2019}, we select 100 class names by a fixed hash rule. Each of six domains contributes disjoint sets of 1,536 target-selection images, 1,024 target-validation images, and 3,072 anchor-pool images from official training data, plus 2,048 official test images. The manifest contains 46,080 images. Class metadata define the benchmark vocabulary but are hidden from screening.

CLIP ViT-B/32 \citep{radford2021} serves as a frozen construction anchor. Within each domain, its first principal-component scores define low, middle, and high strata, each supplying a block of 512 images. For a given target, we exclude its own blocks and enumerate $\binom{15}{3}=455$ combinations, each with 1,536 source images. ResNet-50 and DINOv2-B/14 provide evaluation features \citep{he2016,oquab2024}. All rules share row normalization, a fixed Rademacher projection to 128 dimensions, and post-projection row normalization. Holding sample count constant removes a direct source of training-cost differences between candidates.

Target A-opt uses $P_0^{-1}=I$ and $\sigma^2=1$. Comparators are second-moment MMD, target energy $\tr(\widehat C_TG_S/n_S)$, Bayesian D-opt, merged effective rank, subspace DPP, domain balance, and 20 seeded random rankings. They score the same candidate family. Domain balance additionally reads source-domain identifiers. Kernel definitions and tie rules appear in Appendix~\ref{app:protocol}.

The screen produces nested shortlists of sizes $10,25,50,100,227$, with $227$ designated primary. In the original protocol, U2 independently evaluates each method's primary shortlist using ridge regression with regularization $1$, no intercept, and a fixed softmax temperature. Validation Brier selects the final combination, breaking ties lexicographically. Smaller lists reuse that method's nested evaluations; there is no cross-method cache. The test stage verifies upstream hashes before evaluating every combination to establish an exhaustive audit reference.

The follow-up adds four projection seeds and a multinomial logistic readout with fixed regularization. It keeps the candidate members and data roles unchanged. Screen, validation, and test remain separate processes, but the follow-up reuses each combination's fitted model across methods to audit their shortlists. Its runtime is therefore the cost of a cached research audit rather than separate deployments. Readout settings and provenance checks are detailed in Appendix~\ref{app:protocol}.

We measure how many test top-ten candidates survive screening and how far the validation choice lies from the test oracle $S^*$. Recall measures candidate retention, not classification accuracy; regret measures loss relative to the oracle:
\begin{equation}
 \mathrm{Recall}_{10}(\cQ_s)=\frac{|\cQ_s\cap\mathrm{Top}_{10}(b)|}{10},\qquad
 \mathrm{Regret}(\widehat S)=\frac{b(\widehat S)-b(S^*)}{|b(S^*)|}.
 \label{eq:metrics}
\end{equation}
The original acceptance criteria require at least $50\%$ removal, mean recall of at least $90\%$, and mean normalized regret at most $5\%$. Brier is primary. Follow-up negative log-likelihood (NLL) and classification-error audits use their corresponding validation metric to select candidates, with the metric fixed before test access.

We average encoder results within each dataset--target-domain unit and random rankings within each task. PACS/Office-Home contain eight such units and DomainNet six. Follow-up seeds are also averaged within domains. Original bootstrap intervals resample target units; the test-sample analysis fixes trained models and resamples examples within each target. Seeds, encoders, and overlapping combinations are repeated measurements, not additional independent datasets.

\section{Empirical Findings}
\label{sec:results}
We first examine final choices and computation cost on DomainNet, then ask whether the advantage comes from domain composition or from choosing blocks within a domain. Retention and readout analyses explain which alternatives remain available. The earlier shared-model and direct-selection studies provide context for interpreting these results; they are not additional equal-cost replications.

\subsection{Final choices across evaluation budgets}
\label{sec:budget}
Figure~\ref{fig:budget} reanalyzes the four new projection seeds using final test loss rather than top-ten membership. The reference fits all candidates and selects by the same validation metric. At small budgets, target-aware lists have lower mean Brier than equally sized random lists. Their curves approach the exhaustive-validation reference well before all candidates are evaluated. The practical comparison is therefore between decision procedures at a specified loss tolerance, not simply between ranking correlations.

\begin{figure}[t]
\centering
\includegraphics[width=\linewidth]{decision_budget.pdf}
\caption{Mean final test Brier minus the exhaustive-validation choice is plotted against the number of candidates evaluated. The two panels use different vertical scales. The zero line is a validation-selected model, not the test oracle. Curves are post-hoc summaries of the four new projection seeds; the horizontal axis counts fits, not measured runtime.}
\label{fig:budget}
\end{figure}

For logistic regression, evaluating 50 A-opt or MMD candidates gives lower mean Brier than evaluating 227 random candidates. Both comparisons favor screening in five of six domain averages. The paired mean differences are $-0.001518$ for A-opt and $-0.001664$ for MMD; descriptive domain-bootstrap ranges are $[-0.003517,0.000023]$ and $[-0.003720,-0.000119]$. The domains share source pools, so these are not confidence statements about independent datasets.

The task-level exceedance rates show what the means omit. At size 50, each geometric method exceeds an absolute Brier gap of $0.005$ to exhaustive validation in one of 48 repeated encoder--seed--domain tasks, or $2.08\%$. Random screening at size 227 exceeds it in $15.52\%$ after the prescribed averaging. Appendix Table~\ref{app:budget-table} reports the final losses. This tolerance is one point in an exploratory sensitivity grid, not a new success gate. The comparison describes both average quality and the frequency of larger gaps within the evaluated tasks.

\subsection{Measured decision and feature costs}
\label{sec:cost}
Table~\ref{tab:cost} tests whether evaluating fewer candidates saves time. Each method scores, fits, and validates independently; no fitted model is shared across methods. Quality and time come from the same runs. Unlike the preceding four-seed analysis, this profile fixes one projection and uses three executions per encoder--target--readout task. Appendix Figure~\ref{fig:cost} shows all measured budgets.

With projected features already in memory, A-opt and MMD at size 50 take about one tenth of exhaustive logistic-regression time. Adding the measured image-feature work leaves their accounted cost at about $56\%$ of exhaustive validation. For ridge, extra target-selection features cost more than the cheap fits screening avoids. Thus the same reduction in candidate count has different consequences for the two readouts. Image totals add separately measured components; they exclude pool construction, downloads, cache serialization, and test auditing. Disk-cache loading and projection are reported separately in Appendix Table~\ref{app:cache-table}.

\begin{table}[ht]
\centering
\caption{Matched final Brier and mean accounted time in seconds are reported for the two readouts. Resident time includes scoring, fitting, and validation from in-memory projected features. Images adds the separately measured feature work for the fixed pool inventory; it is an additive profile, not a complete pipeline replay. All means use the same timed choices.}
\label{tab:cost}
\begin{tabular}{llrrrr}
\toprule
Readout & Screen & Fits & Brier & Resident (s) & Images (s)\\
\midrule
Ridge & Target A-opt & 50 & 0.986655 & 0.53 & 39.59 \\
Ridge & MMD & 50 & 0.986643 & 0.21 & 39.27 \\
Ridge & Random & 227 & 0.986653 & 0.75 & 35.30 \\
Ridge & Full validation & 455 & 0.986639 & 1.48 & 36.03 \\
Logistic & Target A-opt & 50 & 0.863170 & 4.49 & 43.54 \\
Logistic & MMD & 50 & 0.863110 & 4.31 & 43.37 \\
Logistic & Random & 227 & 0.864241 & 21.23 & 55.77 \\
Logistic & Full validation & 455 & 0.862992 & 42.83 & 77.38 \\
\bottomrule
\end{tabular}

\end{table}

\subsection{Candidate construction and source-domain controls}
\label{sec:construction}
We next fix the projection and readout parameters and replace the CLIP-stratified blocks with a label- and feature-independent hash partition of the same source images. Mean matching and domain-level moment matching provide simpler controls. The latter pools all source blocks within each domain before scoring domain mixtures, so it cannot distinguish blocks with identical source-domain composition. Appendix Table~\ref{app:construction-table} reports this exploratory challenge; randomized tie handling is reported separately from the frozen lexicographic protocol.

The new construction preserves the broad Brier advantage over random screening, but changes which geometric details help. At size 50, domain-level matching and MMD have similar mean final Brier after repartitioning. Good global shortlists can therefore be obtained here using domain-level moments without distinguishing individual blocks within a domain. Mean matching is competitive on the original blocks but less effective after repartitioning.

To compare blocks while holding domain composition fixed, we group candidates that contain the same number of blocks from each source domain. We retain the floor of half the candidates in each non-singleton group. With the original blocks and logistic readout, A-opt and MMD retain the empirical Brier optimum in about $77\%$ and $76\%$ of groups, compared with $45\%$ for random selection. After repartitioning, these rates fall to about $56\%$ and $55\%$; for classification error, both fall below random retention. These are within-group optimum-retention rates, not global top-ten recall. They show useful block information on the original construction, with a benefit that depends on construction and metric.

\subsection{Retaining promising equal-cost combinations}
The original DomainNet study asks which candidates remain available after screening (Table~\ref{tab:primary}). A-opt retains nearly all test top-ten candidates while discarding approximately half the combinations. Every primary A-opt and MMD shortlist contains a test Brier oracle; their remaining selected regret comes from validation choosing another member. This illustrates Equation~\eqref{eq:decomposition}: retaining the best candidate and selecting it are separate outcomes.

\begin{table}[t]
\centering
\caption{The original DomainNet study retains $s=227$ of 455 combinations. Means average encoders within six target domains. Recall concerns the test Brier top ten; regret is relative to oracle Brier. Spearman measures agreement between the complete screening and test-loss rankings.}
\label{tab:primary}
\begin{tabular}{lrrr}
\toprule
Screen & Recall (\%) & Regret (\%) & Spearman\\
\midrule
Target A-opt & 98.333 & 0.001311 & 0.7443 \\
Second-moment MMD & 98.333 & 0.001311 & 0.6902 \\
Bayesian D-opt & 82.500 & 0.002658 & 0.4996 \\
Merged effective rank & 81.667 & 0.002909 & 0.4951 \\
Random (20 repeats) & 48.958 & 0.003899 & 0.0016 \\
Target energy & 46.667 & 0.009347 & -0.0399 \\
DPP subspace & 27.500 & 0.009067 & -0.0633 \\
Domain balance & 25.000 & 0.012083 & 0.0681 \\
\bottomrule
\end{tabular}

\end{table}

MMD reaches the same primary recall and selected regret as A-opt, and validation chooses the same final combination for both in every encoder--target task. At smaller shortlist sizes, MMD has higher mean recall (Appendix Figure~\ref{fig:recall}). A-opt's stronger full-ranking correlation therefore does not translate into better retention at those operating points. The common advantage over the tested target-blind rules supports using target-conditioned second-order information; direct moment matching is already effective here.


\subsection{Projection and readout sensitivity}
\label{sec:readout}
The original ridge-softmax losses differ little: every candidate lies within the $5\%$ regret gate, so that gate does not distinguish screens. The exploratory follow-up asks whether high recall persists with four new projections and logistic regression, keeping source membership and sample cost fixed. Appendix Table~\ref{tab:readout} reports the new-seed averages; the original seed is an implementation control.


A-opt retains $\ReadoutRidgeA\%$ of the Brier top ten on average under each readout, with MMD slightly higher. The mean relative best-to-worst Brier span increases from $\RidgeSpan\%$ under ridge-softmax to $\LogisticSpan\%$ under logistic regression. High recall thus persists in a setting with more separation between candidates, rather than appearing only when their Brier losses are almost equal. The two target-aware scores also retain most NLL- and error-ranked candidates on average, with MMD higher in these comparisons.

At the primary shortlist size, both methods retain a Brier test oracle in every new-seed task and have the same mean selected regret. Metric-specific outcomes vary: one A-opt task has only $\WorstErrorRecall\%$ error-ranked recall, and validation misses a retained test-error oracle. Appendix~\ref{app:diagnostics} gives the task details and full comparator table.

\subsection{Target weighting under the shared risk model}
The earlier synthetic study tests target weighting when its risk model holds. It varies feature dimension, candidate count, source budget, target sample size, and target rank, using 20 generation seeds for 216 design settings. Under the shared conditional model, A-opt reduces modeled risk by $4.23\%$ relative to the strongest target-blind effective-rank comparator and by $14.85\%$ relative to MMD. The respective descriptive 95\% intervals are $[3.86,4.62]\%$ and $[13.99,15.70]\%$, resampling design settings after averaging seeds. These comparisons concern modeled linear risk, rather than the classification losses above.

A separate source-conditional-shift sweep changes the regression mechanism. Its mean normalized prediction regret increases from $0.211$ without shift to $3.143$ at the largest tested shift. This endpoint measures realized prediction rather than posterior risk. Together, the two studies distinguish the advantage of covering target-relevant directions from the effect of transferring labels across different conditional mechanisms.

\subsection{Direct selection and available improvement headroom}
PACS and Office-Home test whether the score can directly choose a better source combination. At block budget $K=3$, A-opt improves mean within-unit relative Brier by $0.351\%$ over DPP and $0.273\%$ over MMD. Both fall below the prespecified $2\%$ practical-effect threshold. The paired difference interval crosses zero against DPP and is negative against MMD; detailed values are in Appendix~\ref{app:diagnostics}. These comparisons hold block count, rather than source sample count, fixed.

The post-hoc audit evaluates every feasible combination to measure improvement headroom: the loss reduction available by choosing the test-best candidate rather than the reference. Appendix Figure~\ref{fig:diagnostics} shows only $0.781\%$ mean headroom over DPP. Most target units leave little room for a selector to improve the average loss substantially. This distinguishes direct predictive gains from the screening question: whether fewer combinations can be evaluated while retaining useful choices.

\section{Discussion}
\subsection{What makes screening useful?}
The experiments assess screening through three connected outcomes: which candidates survive, which predictor validation selects, and how much work the decision requires. High recall alone answers only the first question. The MMD and domain-matching results also make simple target comparisons essential baselines for more detailed block geometry.

The operating point should reflect the target metric, acceptable loss gap, and measured cost, rather than retention alone. Compression introduces another requirement: the score must still distinguish candidates. In Appendix~\ref{app:diagnostics}, valid sketch intervals remain too wide to resolve choices at a useful byte budget.

\subsection{Limitations}
\label{sec:limitations}
The equal-sample-cost evidence comes from two constructions on the same DomainNet images. Its target domains share source pools, and the added seeds and readouts reuse the same images. They measure sensitivity within this benchmark rather than replication on independent data. Changes from PACS/Office-Home also include the dataset, projection, sample budget, and candidate construction, so they do not isolate candidate heterogeneity as a cause.

The Bayesian derivation concerns shared linear regression; the classification results establish the usefulness of its score empirically. Test-example resampling conditions on fitted models and validation choices. It does not measure uncertainty from training or validation selection. Absolute task quality also depends on the candidate family, as the supplementary adapter study illustrates. Image costs are measured components, not a replay of a complete deployment: pool construction and cache serialization are excluded, and the feature profile uses one pass on one hardware setup. Independent equal-cost datasets and controlled validation budgets remain important, as does checking overlap with encoder pretraining data.

\section{Conclusion}
\label{sec:conclusion}
We evaluated target-aware geometry as a screening step before supervised data-pool selection. The analysis connects experimental design to shared linear risk and separates candidate omission from validation error. On equal-sample-cost DomainNet combinations, evaluating 50 A-opt or MMD candidates gives lower mean test Brier than evaluating 227 random candidates. Domain controls show useful but construction-dependent block information, with MMD competitive in the DomainNet screening comparisons. Independent timing supports savings with cached features; adding measured feature costs preserves the logistic-regression saving but reverses it for ridge. The resulting evidence supports judging a screen by its final choice, acceptable loss gap, and measured cost together.
